\documentclass[conference]{IEEEtran}
\IEEEoverridecommandlockouts
% The preceding line is only needed to identify funding in the first footnote. If that is unneeded, please comment it out.
\usepackage{cite}
\usepackage{amsmath,amssymb,amsfonts}
\usepackage{algorithmic}
\usepackage{graphicx}
\usepackage{textcomp}
\usepackage{xcolor}
\def\BibTeX{{\rm B\kern-.05em{\sc i\kern-.025em b}\kern-.08em
    T\kern-.1667em\lower.7ex\hbox{E}\kern-.125emX}}
\begin{document}

\title{AANEC-Relaxed Trapped Surface Theorem and Stellar Mass-Gap Derivation}\author{Anonymous Research Plan Author}\maketitle

\begin{abstract}
This proposal introduces an AANEC-relaxed trapped surface theorem designed to establish a compactness bound for non-black-hole stellar remnants within a restricted model class of spherically symmetric, regular spacetimes with complete null generators. The planned contribution replaces pointwise energy conditions with an explicit averaged null-convergence assumption, aiming to derive a threshold condition that excludes stable non-black-hole remnants above a model-dependent compactness limit.

The expected outcome is a formal proof or counterexample distinguishing whether the stated assumptions imply a trapped-surface bound, alongside a conditional assessment of how such a bound interacts with astrophysical confounders like explosion energetics and equation-of-state uncertainties. No observed results are claimed; the focus remains on verifying the logical consistency of the theorem and its separation from empirical mass-gap explanations.
\end{abstract}
\begin{IEEEkeywords}
trapped surface, AANEC, compactness, stellar collapse, mass gap, Raychaudhuri equation, Misner-Sharp compactness, null generators
\end{IEEEkeywords}\section{Introduction}
\subsection{Compact Remnants and the Stellar Mass Gap}
The observed stellar mass gap is a central feature in the population of compact remnants, yet its physical origin remains debated. Gravitational-wave observations and stellar-evolution modeling indicate that the location of the lower edge of the black hole mass gap is dictated by the near-final core mass, which in turn sets the resulting black hole mass. This suggests that the gap is astrophysically inferred and governed by progenitor physics rather than being a direct consequence of a pure general-relativistic theorem. Consequently, the gap serves as a consistency constraint for models of massive-star collapse, requiring a clear separation between formal geometric bounds and the complex stellar-physics mechanisms that populate the observed remnant distribution.~\cite{cite_fd0f77b2a173c0ab,cite_7a407c865faafd11}

\subsection{Alternative Explanations and Confounders}
While geometric bounds may restrict the final state of collapse, the observed distribution of remnant masses is heavily influenced by alternative stellar-physics mechanisms. Numerical simulations of binary mergers demonstrate that the use of a non-isentropic equation of state leads to delayed collapse to a black hole, particularly in high-mass binaries, due to pressure support from temperature increases via shocks. Furthermore, stellar-evolution codes incorporating diffusion for angular momentum and chemical abundances enable calculations of rotating-star models, highlighting how rotation and mixing can alter collapse outcomes. These findings indicate that explosion energetics, fallback, equation-of-state uncertainties, and binary interactions act as confounders that can populate or avoid the stellar mass gap independently of a trapped-surface theorem.~\cite{cite_7a44e6730c44cb39,cite_bfae3a92efd1d48f}

\subsection{Proposal Scope and Claim Discipline}
This proposal introduces a formal theorem for classical general relativity under restricted conditions, specifically spherical nonrotating collapse with isotropic pressure. The central contribution is to replace pointwise energy conditions with the Averaged Null Energy Condition (AANEC) combined with an explicit averaged null-convergence assumption on null congruence generators. This formal mechanism is designed to derive a trapped-surface threshold via Raychaudhuri-type focusing and Misner-Sharp compactness. Crucially, the proposal maintains strict claim discipline: the observed 2-5 solar mass gap is not assumed to be a direct consequence of this theorem. Instead, it is treated as an astrophysical consistency constraint, subject to the alternative mechanisms identified in the literature.

\subsection{Design Assumptions and Boundary Conditions}
The validity of the proposed formal bound depends on several design assumptions. The model class assumes globally hyperbolic spacetime with regular matter fields and complete null generators. It further assumes that non-black-hole remnants are modeled as stable equilibrium objects obeying a bounded compactness limit. The theorem is expected to fail or require restriction if AANEC is violated, if the added averaged null-convergence condition fails, or if nonspherical, rotating, or magnetic dynamics dominate the collapse. These boundary conditions define the scope of the formal contribution, separating it from the broader astrophysical inference required to connect the bound to the observed mass gap.

\section{Background, Survey, and Research Gap}
\subsection{Theoretical Foundations and Energy Condition Tensions}
The proposed research addresses the validity of Penrose-type singularity theorems in regimes where classical pointwise energy conditions are violated. Survey evidence indicates that standard energy conditions (NEC, WEC, SEC) are systematically incompatible with quantum field theory, as phenomena such as the Casimir effect and squeezed vacuum states exhibit local negative energy densities [survey:survey\_markdown\#section-003]. This tension challenges the direct application of classical trapped-surface arguments. To preserve the predictive power of singularity theorems, the literature suggests a shift toward averaged conditions, specifically the Achronal Averaged Null Energy Condition (AANEC), which bounds the magnitude and duration of negative energy along null geodesics [survey:survey\_markdown\#section-003]. However, the survey explicitly notes that the rigorous formulation of these bounds for interacting fields remains hindered by renormalization ambiguities, and the derivation of AANEC from first principles in quantum gravity is an open problem.~\cite{cite_0ed84243cf171bc4}

\subsection{The Unresolved Bridge: AANEC and Trapped Surfaces}
A critical gap identified in the current theoretical landscape is the precise logical bridge between AANEC and the formation of trapped surfaces. While AANEC is posited as a robust pathway to preserve singularity theorems, the survey highlights that it may not imply the local or integrated focusing required for Penrose-type results without additional assumptions [survey:survey\_markdown\#section-003]. Specifically, it remains unclear under which validity domains AANEC alone, or AANEC combined with specific global hyperbolicity constraints, forces the convergence of null congruences necessary for trapped-surface formation. This uncertainty represents a missing premise in the formal claim structure for sub-hypothesis SH3, preventing a definitive evaluation of whether relaxed energy conditions suffice to exclude non-black-hole remnants [survey:survey\_markdown\#section-003].

\subsection{Alternative Compact Objects and Boundary Variables}
The exclusion of stable non-black-hole remnants relies on bounding the compactness of exotic alternatives. Survey evidence describes a taxonomy of Exotic Compact Objects (ECOs) that evade classical bounds by violating isotropy or energy conditions, often characterized by a closeness parameter epsilon relative to the Schwarzschild radius [survey:survey\_markdown\#section-004]. These objects can mimic black hole signatures if their surface is sufficiently close to the horizon, but they are subject to instabilities, such as the ergoregion instability, which limits their spin and longevity [survey:survey\_markdown\#section-004]. A significant research gap exists in identifying a universal boundary variable that rigorously separates stable ECOs from black holes within the restricted model class. The survey notes that while tidal Love numbers and gravitational wave echoes offer observational discriminants, their theoretical scaling laws are model-dependent, leaving the precise compactness threshold for stable remnant exclusion unconstrained by current formal theorems [survey:survey\_markdown\#section-004].

\subsection{Astrophysical Inference Chains and Mass Gap Confounders}
The observed stellar mass gap between neutron stars and black holes is frequently cited as evidence for a fundamental compactness limit. However, survey findings demonstrate that this gap is primarily governed by astrophysical mechanisms rather than pure general-relativistic theorems. The lower edge of the mass gap is dictated by near-final core mass and supernova explosion energetics, with rapid instabilities favoring direct collapse or strong explosions that bypass intermediate remnant masses [survey:survey\_markdown\#section-006]. Alternative explanations, including fallback, binary evolution, metallicity-driven wind stripping, and equation-of-state uncertainties, significantly modulate the remnant distribution [survey:survey\_markdown\#section-007]. Consequently, the mass gap serves as an astrophysical consistency constraint rather than a direct proof of a trapped-surface theorem. The inference chain from observed mass distributions to formal geometric bounds is thus confounded by stellar physics, requiring careful separation of formal predictions from observational demographics.~\cite{cite_7a407c865faafd11}

\subsection{Synthesis of Missing Premises and Review Requirements}
The synthesis of formal and empirical evidence reveals two primary unresolved premises. First, the formal claim that AANEC plus explicit averaged focusing implies a trapped-surface bound lacks a rigorous derivation in the restricted model class, as the specific mathematical form of the averaged focusing condition and its compatibility with global hyperbolicity remain undefined [survey:survey\_markdown\#section-003]. Second, the boundary variable that definitively excludes stable non-black-hole remnants above a compactness threshold is not uniquely identified, as ECO stability depends on model-specific parameters not covered by general theorems [survey:survey\_markdown\#section-004]. These gaps necessitate human review to determine whether the proposed theorem should be restricted to classical matter fields or extended with explicit quantum-corrected focusing assumptions. Until these premises are resolved, the derivation of a mass-gap exclusion from AANEC remains a design assumption rather than a verified theoretical consequence.

\section{Research Questions and Planned Contributions}
\subsection{Formalizing the AANEC-Relaxed Trapped Surface Threshold}
The first research question asks whether spherically symmetric stellar collapse, governed by classical general relativity with regular matter fields, global hyperbolicity, and complete null generators, necessarily forms a trapped surface once a specific compactness threshold is crossed. This inquiry requires replacing standard pointwise energy conditions with the Averaged Null Energy Condition (AANEC) supplemented by an explicit averaged null-convergence/focusing condition. A critical unresolved premise is the exact mathematical form of this additional focusing condition, which must be distinguished from standard AANEC to serve as an independent geometric hypothesis. Without this precise definition, the logical bridge between energy condition satisfaction and trapped-surface formation remains formally incomplete.

\subsection{Defining Compactness and Stability for Non-Black-Hole Remnants}
The second research question investigates the exclusion of stable non-black-hole remnants above a compactness-dependent mass scale. This requires a rigorous definition of the Misner-Sharp compactness parameter for the specific metric class under consideration, as well as a formal stability predicate that distinguishes these objects from regular black hole models like Hayward or Bardeen solutions. The proposal posits that if such objects are modeled as stable equilibrium states obeying a bounded compactness or TOV-like limit, they may be excluded by the theorem. However, the current context lacks the formal definitions necessary to operationalize these stability criteria and compactness thresholds, delimiting the scope of the inquiry to a theoretical framework that is yet to be fully specified.

\subsection{Planned Contribution: Separating Geometric Bounds from Astrophysical Explanations}
The primary planned contribution is to establish a qualified AANEC trapped-surface bound that explicitly separates geometric theorem constraints from astrophysical mass-gap explanations. The proposal aims to derive a trapped-surface threshold via Raychaudhuri-type focusing and Misner-Sharp compactness, thereby excluding stable non-black-hole remnants within the restricted model class. Crucially, this contribution does not claim to quantitatively derive the observed 2-5 solar mass gap. Instead, it treats the observed gap as an astrophysical consistency constraint subject to alternative explanations such as explosion energetics, fallback, equation-of-state physics, and binary evolution. The contribution lies in clarifying the boundary conditions under which the geometric theorem fails or must be restricted, particularly when quantum stress-energy effects or nonspherical dynamics dominate.

\subsection{Planned Contribution: Counterexample Search and Verification Strategy}
A secondary planned contribution is the systematic search for counterexamples that satisfy all stated assumptions but fail to form a trapped surface. This involves evaluating candidate witnesses to determine if they genuinely violate the explicit averaged null-convergence condition or if the condition itself is ill-defined in the current context. The proposal includes a potential strategy for numerical verification of the Raychaudhuri equation and trapped surface formation for a discrete subset of test metrics. This verification plan is contingent upon resolving the open design questions regarding the mathematical forms of the key assumptions. The contribution is thus defined by the methodological framework for testing the theorem's validity, rather than by a pre-assumed proof or observed result.

\section{Idea Source Checkpoints and Direction Selection Audit}
\subsection{Scientific Attraction and Initial Scope}
The proposal originates from the 'mechanism\_replacement' direction, which seeks to relax classical pointwise energy conditions in favor of the Averaged Null Energy Condition (AANEC) combined with an explicit averaged null-convergence assumption. The scientific attraction lies in deriving a trapped-surface threshold for spherically symmetric collapse that excludes stable non-black-hole remnants above a specific compactness limit. This formal bound is positioned not as a direct predictor of the observed 2-5 solar mass gap, but as a geometric constraint on the existence of exotic compact objects, treating the astrophysical gap as a consistency check rather than a theorem-derived output.

\subsection{Audit of Checkpoints and Unsupported Links}
Audit snapshots from the idea generation pipeline (checkpoints 1-3) reveal critical gaps in the formalization of the proposed mechanism. While the central hypothesis consistently links AANEC and focusing conditions to trapped-surface formation, the operational definitions for the key variables remain incomplete. Specifically, the mathematical form of the averaged null-convergence assumption (A5) and the Misner-Sharp compactness parameter (D4) are missing. Furthermore, the checkpoints explicitly flag that AANEC alone may not imply the local focusing required for Penrose-type arguments, necessitating an independent geometric hypothesis that has not yet been rigorously defined or proven.

\subsection{Qualified Retention and Explicit Exclusions}
The direction is retained under a qualified design assumption that strictly separates formal geometric proofs from astrophysical inference. The selection process excludes any claim that the theorem quantitatively explains the stellar mass gap, acknowledging instead that explosion energetics, fallback, and binary evolution are dominant confounders. The retained scope is limited to proving whether the added averaged focusing condition, when satisfied, forces trapped-surface formation within the restricted model class of regular, globally hyperbolic spacetimes. All other links to observational data are treated as secondary consistency constraints rather than primary evidentiary support.

\section{Problem Definition, Assumptions, and Hypotheses}
\subsection{Problem Scope and Restricted Model Class}
The proposed research addresses a formal problem in classical General Relativity: determining whether spherically symmetric stellar collapse with isotropic pressure forces trapped-surface formation when a model-dependent compactness threshold is exceeded. The work is restricted to a specific model class defined by global hyperbolicity, regular matter fields, and complete null generators. Within this domain, the project seeks to establish whether stable non-black-hole remnants above a compactness-dependent mass scale can be excluded. This formal scope is explicitly separated from astrophysical inference; the observed 2-5 solar mass gap is treated as an external consistency constraint rather than a direct prediction of the theorem. The proposal does not claim to resolve stellar-evolution mechanisms, supernova engine physics, or binary interactions, but rather to clarify the geometric boundaries within which such astrophysical explanations must operate. The central problem is thus a mathematical conjecture about the interplay between averaged energy conditions and compactness bounds, not a quantitative derivation of observed remnant distributions.

\subsection{Assumption Ledger and Energy Condition Relaxation}
\begin{itemize}
\item The formal argument rests on a specific ledger of declared assumptions that distinguish it from classical pointwise energy-condition theorems. First, the geometry assumes spherically symmetric collapse of massive stars with isotropic pressure, global hyperbolicity, and sufficient regularity for trapped-surface arguments. Second, null generators in the relevant region are assumed complete. Third, matter fields satisfy the Averaged Null Energy Condition (AANEC) along these generators. Crucially, the proposal introduces an independent geometric hypothesis: an explicit averaged null-convergence or focusing condition. This condition is not presented as a consequence of AANEC alone, addressing the concern that AANEC may be satisfied while local or integrated focusing fails. Fourth, non-black-hole remnants are modeled as stable equilibrium objects obeying a bounded compactness or TOV-like limit. The relaxation of pointwise NEC/SEC to AANEC plus explicit averaged focusing constitutes the primary methodological shift. However, the precise mathematical form of the averaged focusing condition remains an unresolved premise, and its relationship to AANEC requires human review to ensure it is not merely a restatement of the original pointwise assumption in integral form.
\end{itemize}

\subsection{Formal Objects and Operational Definitions}
\paragraph{Definition.} The formal objects are defined through a candidate symbol set and variable mapping. The spacetime metric is denoted by g, and the stress-energy tensor by T. The mass of the collapsing matter or remnant object is M, while r represents the areal radius or radial coordinate. The expansion scalar of a null geodesic congruence is theta, and k denotes the tangent vector to null geodesic generators. The central geometric quantity is the Misner-Sharp compactness parameter, denoted C. This parameter is intended to measure the effective mass within radius r relative to the areal radius, serving as the threshold variable for trapped-surface formation. The proposal treats C as the independent variable governing the compactness bound, while trapped\_surface\_existence and stable\_remnant\_exclusion are the dependent formal outcomes. The operationalization of C, however, is not fully specified. The definition requires a rigorous formula linking M, r, and the metric components g in the chosen spherical ansatz. Without this explicit definition, the compactness threshold remains a candidate formalization rather than a verified mathematical object.

\[
C \approx \frac{2 M}{r}
\]

\subsection{Candidate Hypotheses and Unresolved Status}
\paragraph{Proposition (proposed).} Two candidate propositions structure the proposed derivation. P1 posits that in spherically symmetric stellar collapse satisfying AANEC, global hyperbolicity, regular matter fields, complete null generators, and an explicit averaged null-convergence condition, reaching a model-dependent compactness threshold forces trapped-surface formation. P2 posits that stable non-black-hole remnants above a compactness-dependent mass scale are excluded within this restricted model class. Both propositions are currently status: candidate\_formalization and remain unverified. The forward derivation from P1 to P2 depends on proof obligations PO1 through PO3, which are unresolved. Specifically, the argument requires demonstrating that the averaged focusing condition forces theta to become negative, that negative theta implies a trapped surface when C exceeds a threshold, and that a trapped surface is incompatible with the stable equilibrium limit. These steps are proposed but not proved. The hypothesis mapping identifies H1 and H2 as testable formal claims, but the counterexample analysis remains a design assumption. The final canonical field status for the assumption ledger requires human input, meaning the precise formulation of the averaged focusing condition and the Misner-Sharp compactness parameter must be resolved before the propositions can be considered well-posed.

\section{Study Design and Methods}
\subsection{Analytical Framework and Unit of Analysis}
The study adopts a theoretical and mathematical design, treating formal proof construction and counterexample search as the primary unit of analysis. The proposed method replaces pointwise energy conditions with the averaged null energy condition (AANEC) combined with an explicit averaged null-convergence assumption. This framework aims to derive a trapped-surface threshold using Raychaudhuri-type focusing and Misner-Sharp compactness, subsequently bounding stable non-black-hole remnants via TOV-like stellar-structure assumptions.

\subsection{Operational Definitions and Unresolved Variables}
The operationalization of key variables remains incomplete. Specifically, the mathematical form of the explicit averaged null-convergence condition and the formula for the Misner-Sharp compactness parameter are currently undefined. These definitions are prerequisites for the proposed derivation and require human specification before the formal proof strategy can be fully implemented.

\subsection{Boundary Conditions and Counterexample Strategy}
The validity of the theorem is restricted to spherically symmetric collapse with isotropic pressure, global hyperbolicity, and complete null generators. The design includes a strategy to identify counterexamples that satisfy all assumptions but fail to form a trapped surface. This analysis explicitly excludes astrophysical confounders such as rotation, magnetic fields, and binary evolution from the formal proof, treating them instead as boundary conditions that may restrict the theorem's applicability.

\subsection{Data Governance and Reproducibility}
The protocols for data management and reproducibility in this theoretical context are not yet established. Human input is required to define the standards for documenting the formal derivations, storing intermediate mathematical results, and ensuring the auditability of the proof steps.

\subsection{Measurement and Quality Control}
In the absence of empirical data collection, the measurement plan and quality control procedures refer to the validation of logical consistency and mathematical rigor. These procedural details require human review to define acceptable standards for proof verification and error checking within the formal model.

\section{Expected Outcome Branches and Conditional Conclusions}
\subsection{Supportive Conditional Branch}
If the prespecified analysis is consistent with the declared relation while planned controls do not favor a stated alternative explanation, the result would support, but not prove, the declared relation within the design boundary. This outcome corresponds to the supports\_mechanism branch. The theorem fails or must be restricted if AANEC is violated, the added averaged null-convergence/focusing condition fails, null generators are incomplete, quantum stress-energy violates ANEC, nonspherical/rotating/magnetic dynamics dominate, or stable exotic compact objects evade the averaged-focusing assumptions. The mass-gap inference fails or becomes non-discriminating if explosion energetics, fallback, equation-of-state physics, pair-instability processes, or binary evolution dominate the observed remnant distribution. In this scenario, improvement actions include replicating under independently confirmed conditions and testing the most consequential declared boundary condition.

\subsection{Heterogeneous and Null Branches}
The relation may be conditional or heterogeneous; no universal conclusion is warranted if the prespecified analysis indicates variation across declared conditions, units, or measurement contexts. This corresponds to the partial\_or\_heterogeneous branch, where improvement actions involve predefining and checking plausible moderators and improving the coverage of conditions and measurement comparability. Conversely, if the prespecified comparison does not support the declared relation or instead favors a declared alternative explanation, the proposed relation is not supported in this design boundary. This null\_or\_contradictory outcome indicates that absence of support is not proof of absence generally, requiring an audit of construct validity and comparison adequacy, and revision of the mechanism or boundary claim before another design iteration.

\subsection{Uninformative Branch Consequence}
No scientific conclusion is warranted because the planned design did not yield interpretable evidence if prespecified quality-control, missingness, protocol-deviation, or validity criteria prevent interpretation. This corresponds to the uninformative\_or\_invalid branch. The consequence of this outcome is that the identified validity or data-quality failure must be resolved before repeating the design. Furthermore, human confirmation of measurement, sampling, and analysis prerequisites is required to proceed, as the current design boundary cannot support a definitive conclusion regarding the formal theorem or its astrophysical implications.

\section{Risks, Limitations, and Human Review Requirements}
\subsection{Missing Formal Definitions}
The proposed mechanism relies on two undefined formal objects. The Misner-Sharp compactness parameter C lacks a precise mathematical definition, preventing calculation of thresholds, and the explicit averaged null-convergence/focusing condition A5 has no specified form, making verification of this assumption impossible. These gaps prevent any determination of whether candidate witnesses satisfy all assumptions.

\subsection{Threats from Alternative Explanations}
The observed stellar mass gap may arise from supernova explosion energetics, fallback, equation-of-state uncertainties, rotation, magnetic fields, nonspherical dynamics, or binary interactions, rather than from a trapped-surface theorem. Consequently, the mass gap is treated as an astrophysical consistency constraint, not a direct consequence of the formal claim. The theorem's validity domain is restricted to spherically symmetric, isotropic, regular spacetimes with complete null generators and AANEC satisfaction; violations of these conditions would require restricting or rejecting the bound.

\subsection{Counterexample Verification Status}
The current counterexample analysis is unverified. Candidate witnesses for a stable non-black-hole remnant above the compactness threshold have been proposed, but their validity depends on resolving the missing definitions of A5 and D4. Until these formal gaps are closed, no conclusion can be drawn about whether the theorem is falsifiable within its stated assumptions.

\subsection{Required Human Review Decisions}
Before escalation, human review must resolve the formal reasoning plan status, specify the mathematical form of A5 and D4, and distinguish formal counterexamples from empirical alternatives. This review is required to assess whether the proposed witnesses satisfy all assumptions and to determine if the theorem's scope must be restricted.

\section{Definitions, Propositions, and Proof Obligations}
\subsection{Geometric and Matter Definitions}
The formal framework relies on classical General Relativity in a spherically symmetric setting. The spacetime metric g (D6) describes the geometry of the collapsing matter. The stress-energy tensor T (D5) encodes the matter distribution, assumed to have isotropic pressure. Null geodesic generators are characterized by their tangent vector k (D7), which defines the direction of light propagation in the relevant regions. The expansion scalar theta (D3) measures the rate of change of the cross-sectional area of a null congruence, serving as a key indicator of focusing. The mass M (D1) represents the total mass of the collapsing object or remnant, while the areal radius r (D2) provides the radial coordinate for defining compactness. These definitions establish the variables used in subsequent propositions.

\subsection{Unresolved Compactness Parameter}
A critical component of the argument is the Misner-Sharp compactness parameter C (D4). This quantity is intended to measure the degree of compactness of the object relative to its radius, which determines whether a trapped surface forms. However, the specific mathematical formula for C, including the definition of effective mass within radius r in the chosen metric, is currently missing from the canonical context. Without this precise definition, the threshold for trapped-surface formation cannot be rigorously calculated. Consequently, the definition of C remains a design assumption requiring human input to specify the exact functional form.

\subsection{Core Propositions}
Two central propositions structure the theoretical argument. Proposition P1 posits that in spherically symmetric stellar collapse satisfying the Averaged Null Energy Condition (AANEC), global hyperbolicity, regular matter fields, complete null generators, and an explicit averaged null-convergence/focusing condition, reaching a model-dependent compactness threshold forces trapped-surface formation. Proposition P2 extends this to conclude that stable non-black-hole remnants above a compactness-dependent mass scale are excluded within this restricted model class. These propositions are currently candidate formalizations; their validity depends on the resolution of the missing definitions and the proof of the associated obligations.

\subsection{Proof Obligations and Missing Assumptions}
The derivation of P1 and P2 entails several proof obligations. PO1 requires verifying that the explicit averaged null-convergence condition (A5) holds along the relevant null generators. However, the mathematical form of A5 is not provided, distinguishing it from the standard AANEC (A2). PO2 involves ensuring the global hyperbolicity and regularity of the spacetime metric (A3) are sufficient for trapped-surface arguments. PO3 demands the calculation of the compactness threshold using the defined parameter C. These obligations are unresolved due to the missing formal inputs for A5 and D4.

\subsection{Forward Derivation Status}
The proposed forward derivation proceeds in four steps. Step S1 suggests that the expansion scalar theta becomes negative along null generators due to the averaged focusing condition. Step S2 proposes that if the compactness parameter C exceeds a threshold, both outgoing and ingoing null expansions become negative, indicating a trapped surface. Step S3 checks the consistency of trapped-surface formation against the assumption that non-black-hole remnants are stable equilibrium objects (A6). Step S4 concludes that stable non-black-hole remnants are excluded for mass/compactness regimes where C exceeds the threshold. All steps are currently marked as proposed or unverified, contingent on the resolution of the underlying definitions and assumptions.

\subsection{Scope of Formal Claims}
The template details for the formal claim and proof obligations are currently flagged as requiring human input. This indicates that the precise logical structure and the specific requirements for closing the proof obligations have not been fully specified in the canonical context. The final status of these template fields must be resolved to ensure that the formal argument is complete and that the proposed derivation aligns with the intended theoretical scope.

\section{Forward Derivation and Counterexample Search Plan}
\subsection{Lemma-by-Lemma Forward Derivation Plan}
The proposed forward derivation proceeds through a sequence of four lemmas that link energy conditions to compactness bounds. Step S1 applies a Raychaudhuri-type focusing argument to show that the expansion scalar theta becomes negative along null generators under the explicit averaged focusing condition A5. Step S2 proposes that when the Misner-Sharp compactness parameter C exceeds a model-dependent threshold, both outgoing and ingoing null expansions become negative, indicating trapped-surface formation. Step S3 checks consistency between trapped-surface formation and the stability predicate A6, which models non-black-hole remnants as stable equilibrium objects obeying a bounded compactness or TOV-like limit. Step S4 derives the exclusion bound for stable non-black-hole remnants above a compactness-dependent mass scale. Each step is currently marked proposed or unverified and depends on formal obligations PO1, PO2, and PO3 that remain unresolved.

\[
\int_{\gamma} T_{ab} k^a k^b \, d\lambda \geq 0
\]

\subsection{Admissible Comparators and Failure Conditions}
The theorem's validity is bounded by several admissible comparators and failure conditions. The result fails or must be restricted if AANEC is violated, if the added averaged null-convergence/focusing condition A5 fails, if null generators are incomplete, or if quantum stress-energy effects violate ANEC. Nonspherical, rotating, or magnetically dominated dynamics may also invalidate the spherical isotropy assumption A1. Stable exotic compact objects that evade the averaged-focusing assumptions constitute another boundary case. For the astrophysical mass-gap application, the inference becomes non-discriminating if explosion energetics, fallback, equation-of-state physics, pair-instability processes, or binary evolution dominate the observed remnant distribution. These failure conditions define the scope within which the formal claims are proposed to hold.

\subsection{Valid Counterexamples Versus Out-of-Domain Boundary Cases}
A valid counterexample must satisfy all assumptions A1 through A6 while negating proposition P2. Candidate CE1 proposes a static, spherically symmetric, regular metric with high compactness that satisfies AANEC but potentially violates the explicit averaged focusing condition A5 or the stability bound A6. Candidate CE2 considers regular black hole models such as Hayward or Bardeen metrics that satisfy AANEC and global hyperbolicity but may fail to qualify as non-black-hole remnants if they contain horizons. Both candidates are currently classified as boundary cases or assumptions-not-satisfied rather than valid counterexamples. Distinguishing formal counterexamples from out-of-domain boundary cases requires human review of the specific mathematical structure of the proposed witnesses.

\subsection{Unresolved Formal Obligations and Human Review Needs}
The counterexample and boundary analysis, together with the verification plan, remain fields that require human input. The specific formal definition of the Misner-Sharp compactness parameter C is missing, preventing precise calculation of thresholds. The explicit mathematical form of the averaged null-convergence condition A5 is also absent, making verification of this assumption impossible. These unresolved items signal that the current derivation and counterexample search are provisional and cannot be treated as verified results.

\appendices
\section{Idea Source Checkpoints and Direction Selection Audit}
\subsection{Available Audit Snapshots}
Three available audit snapshots were recovered from the idea generation pipeline: an initial candidate record, a portfolio entry with a qualified title variant, and a direction-specific result record. These checkpoints are presented as available source material rather than as a proven temporal iteration sequence. Each snapshot preserves the core proposal for an AANEC-relaxed trapped-surface theorem and a compactness bound for non-black-hole stellar remnants, while the portfolio entry displays a narrower title that foregrounds the qualified nature of the bound. The direction result record additionally contains explicit lists of assumptions, boundary or failure conditions, discriminating observations, evidence requirements, and expected mechanisms that are not fully duplicated in the other two snapshots.

\subsection{Scope Corrections and Retained Decisions}
The retained direction selection is mechanism replacement, and all three snapshots consistently record the same central hypothesis and claim scope. The proposal replaces pointwise NEC/SEC assumptions with AANEC plus an explicit averaged null-convergence/focusing condition, derives a trapped-surface threshold via Raychaudhuri-type focusing and Misner-Sharp compactness, and bounds stable non-black-hole remnants through TOV-like stellar-structure assumptions. A key scope correction preserved across the checkpoints is that the observed 2-5 solar mass gap is not claimed to be a direct consequence of the theorem; it is treated instead as an astrophysical consistency constraint subject to alternative explanations such as explosion energetics, fallback, equation-of-state uncertainties, rotation, magnetic fields, nonspherical dynamics, and binary interactions.

\subsection{Open Questions and Missing Formal Material}
The audit snapshots identify three target gaps that remain unresolved: a missing formal claim assumption, an unmapped gap boundary variable, and an unresolved progenitor explosion energetics mechanism. The direction result record also flags uncertainty about whether AANEC alone implies the local or integrated focusing needed for Penrose-type trapped-surface formation, making the added averaged null-convergence/focusing condition essential but potentially as restrictive as the original pointwise assumption in some regimes. These open questions are preserved for human review rather than resolved here; no proof, counterexample, or empirical result is asserted.

\section{Variables, Symbols, and Operational Definitions}
\subsection{Geometric and Stress-Energy Symbols}
The formal model defines the spacetime metric as \$g\$ and the stress-energy tensor as \$T\$. The expansion scalar of a null geodesic congruence is denoted by \$\textbackslash{}theta\$, while the tangent vector to the null geodesic generators is represented by \$k\$. These symbols form the basis for the proposed Raychaudhuri-type focusing arguments.

\subsection{Compactness and Mass Definitions}
The mass of the collapsing matter or remnant object is defined as \$M\$, and the areal radius or radial coordinate is defined as \$r\$. The Misner-Sharp compactness parameter, denoted by \$C\$, is intended to relate these quantities. However, the specific formula for \$C\$ and the effective mass within radius \$r\$ are not currently defined, creating an unresolved operational gap.

\subsection{Unresolved Operationalization of Collapse Parameters}
\begin{itemize}
\item The following variables require human input to establish their operational definitions before the formal derivation can proceed:
\item The specific equation of state models or nuclear physics parameters to be tested are not listed.
\item Measurement or initial condition specification for rotation in collapse models is not defined.
\item Measurement or initial condition specification for magnetic fields in collapse models is not defined.
\item The specific formula for Misner-Sharp compactness threshold dependent on model parameters is not defined in the brief.
\item Measurement method for stellar remnant mass (e.g., radial velocity curves vs pulsar timing) is not specified.
\item Specific metric for explosion energy (kinetic vs total) and its measurement in observed samples is not defined.
\end{itemize}

\section{Evidence Coverage, Unknown Items, and Review Checklist}
\subsection{Bounded Evidence Coverage}
The proposal's reliance on averaged energy conditions is bounded by the finding that pointwise energy conditions are systematically violated by quantum fields, which necessitates weaker statements such as quantum energy inequalities or averaged energy conditions for broader validity. This supports the design choice to replace pointwise NEC or SEC assumptions with AANEC, but it does not prove that AANEC alone implies the additional explicit averaged null-convergence condition required by the theorem. Astrophysical context is supplied by the observation that the lower edge of the black hole mass gap is dictated by near-final core mass and is robust against variations in metallicity, internal mixing, and stellar wind mass loss. This anchors the mass-gap discussion in stellar-physics mechanisms rather than treating the gap as a direct prediction of the formal theorem.~\cite{cite_0ed84243cf171bc4,cite_7a407c865faafd11}

\subsection{Formal Unknowns and Counterexample Limits}
Two formal gaps prevent release of any theorem claim. The Misner-Sharp compactness parameter C lacks a supplied mathematical definition, so compactness thresholds and mass scales cannot be computed. The explicit averaged null-convergence condition A5 likewise lacks a formal mathematical form, so candidate witnesses cannot be tested against it. Counterexample analysis therefore remains a boundary case rather than a completed result: a proposed witness is a static, spherically symmetric, regular metric with isotropic pressure and high compactness satisfying AANEC, but its status under A5 and under the stability predicate is unknown. Because the counterexample must satisfy all assumptions, its validity cannot be decided until A5 and C are formally specified.

\subsection{Release Criteria for Future Claims}
\begin{itemize}
\item Before any formal claim is released, human review must confirm the following: first, supply the mathematical definition of Misner-Sharp compactness C and its relation to the model-dependent threshold; second, supply the mathematical form of the explicit averaged null-convergence condition A5 and verify whether it is independent of AANEC; third, audit the proposed counterexample witnesses to determine whether they satisfy A5 and the stability predicate or instead violate an assumption; fourth, confirm that all empirical mass-gap statements remain separated from the formal theorem and are treated only as astrophysical consistency constraints.
\end{itemize}\section*{References}
\nocite{cite_07d6df5d6978a1fc,cite_0ed84243cf171bc4,cite_1e54e3c2b35110ad,cite_2576833bf7ef143a,cite_283d00782f1f9ecb,cite_357dd780a9e3a378,cite_364386a3569b6e66,cite_37ecbded9d629353,cite_42e341ff6471ce89,cite_572fec0a1e5d74d9,cite_5ea4dc0e1b288757,cite_63476e0b002da386,cite_72d7ff391b29c495,cite_789b47fb0241d9ae,cite_7a407c865faafd11,cite_7a44e6730c44cb39,cite_7ae92f06684c60ef,cite_977568cf7b8a55bf,cite_98c5e596c0c34b68,cite_bfae3a92efd1d48f,cite_c369df652181e8f7,cite_cbbced417410bc6c,cite_cc5bc07a219047b4,cite_d3016ce6d1faa8ec,cite_e228b5fd31e37646,cite_e7b42f305bc08498,cite_ecea60a2300e4892,cite_f634820e22d7a238,cite_fd0f77b2a173c0ab}
\bibliographystyle{IEEEtran}
\bibliography{references}
\end{document}
